\subsection{Sharp per-mode certificates and finite-sample interpretation}
\label{app:theory-survival-certificates}

The energy bounds above can be sharpened by using probability normalization.
We give self-contained applications of relative-entropy data processing and
Pinsker's inequality \citep[Theorems~9 and~31]{vanerven2014renyi}; binary-KL inversion is also a
standard statistical-learning tool \citep{garivier2011klucb}. These certificates
concern a specified probability distribution. Applying one at a checkpoint
requires a valid loss bound there; applying it throughout training additionally
requires a uniform energy bound. Neither application establishes descent for
the implemented neural optimizer. All logarithms are natural.

\subsubsection{A sharp certificate from normalized replay cross entropy}

Write $H(w)=-\sum_b w_b\log w_b$, $h_2(a)=H((a,1-a))$, and
\[
 \operatorname{kl}(a,r)
 =a\log\frac{a}{r}+(1-a)\log\frac{1-a}{1-r},
\]
with the usual endpoint conventions. For $0<a<1$ and $D\ge0$, let
$\ell_-(a,D)$ be the unique $r\in(0,a]$ satisfying
$\operatorname{kl}(a,r)=D$; set $\ell_-(1,D)=e^{-D}$.

\begin{theorem}[Sharp normalized replay certificate]
\label{ext:sharp-replay-certificate}
Let $\mathcal B$ be a finite nonempty bank, let $w_b>0$ with
$\sum_{b\in\mathcal B}w_b=1$, and let $p$ be a probability distribution on
an alphabet containing $\mathcal B$. Suppose
\[
 R_w(p)=-\sum_{b\in\mathcal B}w_b\log p_b\le C<\infty,
 \qquad D=C-H(w).
\]
Then $D\ge0$ and, for every $b\in\mathcal B$,
\begin{equation}
 \label{ext:binary-kl-retention}
 p_b\ge\ell_-(w_b,D)>0.
\end{equation}
This coordinate bound is sharp on the closed probability simplex whenever
an outcome other than $b$ is available. Moreover,
\[
 p(\mathcal B)\ge e^{-D},\qquad
 p_b\ge\max\{0,w_b-\sqrt{D/2}\}.
\]
For $w_b<1$, the positive bound
$p_b\ge\exp[-(D+h_2(w_b))/w_b]$ is also valid.
\end{theorem}
\begin{proof}
Extend $w$ by zero off the bank, obtaining $\bar w$. Then
\[
 0\le\operatorname{KL}(\bar w\|p)=R_w(p)-H(w)\le D.
\]
For $a=w_b$ and $r=p_b$, the log-sum inequality on the complementary
outcomes gives
\[
 \sum_{j\ne b}\bar w_j\log\frac{\bar w_j}{p_j}
 \ge(1-a)\log\frac{1-a}{1-r},
\]
so $\operatorname{kl}(a,r)\le D$. For $0<a<1$, its derivative in $r$ is
$(r-a)/(r(1-r))<0$ on $(0,a)$; it diverges at zero and vanishes at $r=a$.
Thus $r<\ell_-(a,D)$ is impossible. Equality is attained by choosing
$r=\ell_-(a,D)$, $p_j=(1-r)w_j/(1-a)$ on the remaining banked outcomes,
and zero mass off the bank. For $a=1$, the inequality is
$-\log p_b\le D$, with equality obtained by placing the remaining mass on
another outcome. If the entire alphabet has one outcome, necessarily
$p_b=1$. For $D>0$, equality can be approached with positive off-bank
probabilities by increasing $r$ slightly and sending off-bank mass to zero. At $D=0$, feasibility
instead forces $p=\bar w$ exactly.

Let $s=p(\mathcal B)$ and $\widetilde p_b=p_b/s$. The decomposition
\[
 \operatorname{KL}(\bar w\|p)
 =-\log s+\operatorname{KL}(w\|\widetilde p)\ge-\log s
\]
proves the total-mass bound. For fixed $r\in(0,1)$,
$\partial_a^2\operatorname{kl}(a,r)=1/[a(1-a)]\ge4$, whereas its value
and first derivative vanish at $a=r$. Hence
$\operatorname{kl}(a,r)\ge2(a-r)^2$, including endpoints by continuity,
which proves the Pinsker bound. Finally,
\[
 \operatorname{kl}(a,r)
 =-h_2(a)-a\log r-(1-a)\log(1-r)
 \ge-h_2(a)-a\log r
\]
Combining this inequality with $\operatorname{kl}(a,r)\le D$ yields the
positive exponential coordinate bound.
\end{proof}

For the categorical potential argument behind
Eq.~\eqref{eq:replay-retention-bound} and its positive-weight extension,
one may substitute
\[
 C=R_w(p(T))+
 \frac{\Psi(1)-\Psi(P(T))}{\rho_w}.
\]
Normalization improves the termwise certificate $p_b\ge e^{-C/w_b}$.
As a hypothetical example, a uniform bank of 16 with
$C=\log16+0.01$ yields $p_b\ge0.0339712$, while the termwise bound is
approximately $4.62\times10^{-20}$. These are illustrative mathematical
inputs, not measured replay losses. A large excess cross entropy can still
make even the sharp bound too small to be useful.

\subsubsection{What a scalar entropy value can certify}

\begin{theorem}[Sharp entropy threshold for all-mode retention]
\label{ext:entropy-retention-threshold}
Let $q$ be a distribution on a known alphabet of $m\ge2$ possible modes,
allowing zero coordinates. A scalar bound $H(q)\ge h_0$, with
$0\le h_0\le\log m$, implies a strictly positive uniform coordinate floor
if and only if $h_0>\log(m-1)$. In that case the sharp floor is
$q_b\ge r_m(h_0)$ for every $b$, where $r_m(h_0)\in(0,1/m]$ uniquely solves
\[
 h_2(r)+(1-r)\log(m-1)=h_0,
\]
where $m$ counts all possible modes, including those absent from the
evaluation sample.
\end{theorem}
\begin{proof}
Fix $r=q_b<1$. Splitting off that coordinate gives
\[
 H(q)=h_2(r)+(1-r)H\!\left(\frac{q_{-b}}{1-r}\right)
 \le h_2(r)+(1-r)\log(m-1)=:f_m(r).
\]
Uniform remaining coordinates attain equality. On $0<r<1/m$,
$f_m'(r)=\log[(1-r)/(r(m-1))]>0$, with
$f_m(0)=\log(m-1)$ and $f_m(1/m)=\log m$. Thus a coordinate below the
stated root contradicts $H(q)\ge h_0$; a coordinate equal to one already
exceeds the root. The distribution with one coordinate at the root and the
others uniform proves sharpness. Conversely, if $h_0\le\log(m-1)$, one
zero coordinate with the rest uniform satisfies the entropy bound. Even
restricting to strictly positive coordinates permits arbitrarily small
values by approximation.
\end{proof}

For uniform $u$, $H(q)=\log m-\operatorname{KL}(q\|u)$, and data processing
gives $\operatorname{kl}(q_b,1/m)\le\operatorname{KL}(q\|u)$. A forward-KL
budget therefore excludes every missing coordinate precisely when it is
below $\log[m/(m-1)]$. At the threshold, one missing coordinate is possible,
whereas $\operatorname{KL}(u\|q)=\infty$. This compares level sets; it does
not assert that an exact-entropy gradient flow loses modes. Conditional
entropy alone also leaves total correct mass unconstrained, and entropy
estimated from a finite sample is not the exact $H(q)$ of a known alphabet.

\subsubsection{From complete-response likelihood to an execution key}

\begin{corollary}[Response-to-key certificate]
\label{ext:response-key-certificate}
Fix one prompt $x$, one checkpoint, and distinct complete responses $e_j$
with verified keys $b_j$. Give them weights $w_j>0$ summing to one, and put
$W_b=\sum_{j:b_j=b}w_j$. Under a fixed decoding law, suppose
\[
 -\sum_j w_j\log\pi(e_j\mid x)\le C<\infty.
\]
For every represented key $b$,
\[
 \operatorname{kl}(W_b,p(b\mid x))\le C-H(w),\qquad
 p(b\mid x)\ge\ell_-(W_b,C-H(w))>0.
\]
Individually, $p(b_j\mid x)\ge\pi(e_j\mid x)$.
\end{corollary}
\begin{proof}
The distinct complete responses are disjoint events. Extend their target
weights by zero to the full response alphabet; include invalid outcomes
and, if necessary, nontermination as additional outcomes so that the policy
law is normalized. The resulting response-level relative entropy is
$-\sum_jw_j\log\pi(e_j\mid x)-H(w)\le C-H(w)$. Apply the log-sum inequality
within each execution-key event and then within its binary partition into
key $b$ and all other outcomes. The target assigns mass $W_b$ to the first
event and the policy assigns $p(b\mid x)$, so the resulting divergence is
$\operatorname{kl}(W_b,p(b\mid x))$. The inverse argument of
Theorem~\ref{ext:sharp-replay-certificate} proves the floor. Finally, each
complete-response event is contained in its verified-key event, proving
the individual bound used in Lemma~\ref{lem:shared-exemplar-retention}.
\end{proof}

The probability being scored must include EOS or another specified
termination event, or a fixed deterministic horizon. It must use the same
prompt formatting, temperature, vocabulary mask, top-$p$/top-$k$ rule,
response-truncation rule, and verifier as the probability being bounded.
A prefix likelihood is insufficient because its continuations can produce
other keys or invalid outputs. All likelihoods entering one cross entropy
must use the same checkpoint. For a complete sequence of length $L_j$ and
mean-token log score $\bar\ell_j$, its probability is
$\exp(L_j\bar\ell_j)$, not $\exp(\bar\ell_j)$.

Length-normalized loss requires normalization before the theorem applies:
with $a_j=\alpha_j/L_j$, $A=\sum_j a_j$, and $w_j=a_j/A$,
\[
 \sum_j\alpha_j\frac{-\log\pi(e_j\mid x)}{L_j}=A R_w.
\]
Thus an upper bound $D_0$ on this loss supplies $C=D_0/A$; uniform weights
on mean-token scores generally give inverse-length target weights on
complete responses. Different prompts require separate conditional
calculations or an explicitly normalized joint prompt mixture. A change
in one exemplar's likelihood need not equal the change in its whole key's
probability. Approximate log scores need controlled numerical error to
supply rigorous numerical certificates.

\subsubsection{Finite observations and repeated checkpoint measurement}

At a frozen checkpoint, for a key specified before sampling, $N$ independent
draws under the same decoding law give
$X\sim\operatorname{Binomial}(N,p_b)$. Exact binomial-tail inversion gives
one-sided confidence bounds \citep{clopper1934binomial}. If $X=0$, the
one-sided $(1-\alpha)$ upper bound and lower bound are respectively
\[
 U_0=1-\alpha^{1/N},\qquad 0.
\]
Indeed, $\Pr_{p_b}(X=0)=(1-p_b)^N<\alpha$ whenever $p_b>U_0$; for nonzero
counts, the full confidence procedure uses the corresponding binomial-tail
inversion. As hypothetical examples, zero sightings in 32 draws give an
upper bound of $0.08937$ at the one-sided 95\% level, or $0.16495$ when
Bonferroni protects 16 prespecified keys simultaneously. Two sightings in
32 draws instead give a one-sided 95\% lower bound of about $0.01122$ for
one prespecified key. The upper and lower examples are separate one-sided
limits; a joint interval requires allocation to both tails. Zero sightings
do not establish exact extinction.

For repeated inspection of draws from a fixed policy, a confidence sequence
can provide validity under optional stopping \citep{howard2021confidence}.
For fixed beta-mixture parameters $a_0,b_0>0$, let $S_t$ be the number of
successes in $t$ independent Bernoulli draws and let $\mathrm B$ denote
the beta function. Define
\[
 M_t(p)=\frac{\mathrm B(S_t+a_0,t-S_t+b_0)}
 {\mathrm B(a_0,b_0)p^{S_t}(1-p)^{t-S_t}},\qquad
 \mathcal I_t=\{p\in(0,1):M_t(p)<1/\alpha\}.
\]
For each true $p\in(0,1)$, the numerator integrates the likelihood under a
fixed beta distribution. Hence $M_t(p)$ is a mixture of likelihood-ratio
martingales, is nonnegative, and has initial value one. The nonnegative
martingale maximal inequality therefore gives
$\Pr_p(\exists t:p\notin\mathcal I_t)\le\alpha$. Error allocations across
a fixed roster of keys make coverage simultaneous by a union bound;
independence across keys is unnecessary.

This construction requires a fixed Bernoulli probability. Counts from
different prompts, seeds, or changing checkpoints cannot be pooled to
certify the current checkpoint's probability. Separate frozen-checkpoint
measurements can instead allocate error across prespecified
prompt--key--checkpoint statements. Repeated deterministic teacher-forced
scores are not independent Bernoulli observations, and choosing identities
after inspecting the same samples also needs selection-aware inference.
A probability bound at observed checkpoints alone makes no claim about
unobserved intervening or subsequent training states.

For an unknown catalogue, missing mass is the total probability of categories
not yet observed under a fixed iid law. Concentration results address this
quantity \citep{mcallester2003missing,berend2013missing}; it is neither the
number of missing modes nor proof of their extinction. Such bounds do not
certify each named key, and stationary missing-mass formulas do not directly
apply to a history generated by changing policies. Named-key binomial counts
provide the direct measurement of per-mode visibility.
